\begin{titlepage}
\begin{spacing}{1.15}
\bfseries
\vspace*{\fill}
\begin{center}
{\fontsize{16pt}{19.2pt}\selectfont\MakeUppercase{\thesistitle}\par}
\vfill
SKRIPSI\par
\programstudy\par
\vspace{1cm}
\includegraphics[width=4cm]{logo-isi.png}\par
\vspace{1cm}
Oleh:\\
\researchername\\
NIM: \researchernim\par
\vfill
Tugas Akhir Program Sarjana\par
\academicyear\par
\vfill
\MakeUppercase{\programstudy}\\
JURUSAN \MakeUppercase{\departmentname} \MakeUppercase{\facultyname}\\
\MakeUppercase{\institutionname}\\
\graduationyear\par
\end{center}
\vspace*{\fill}
\end{spacing}
\end{titlepage}

\setcounter{page}{2}

% Submission Page
\newpage
\begin{spacing}{1.15}
\bfseries
\vspace*{\fill}
\begin{center}
{\fontsize{16pt}{19.2pt}\selectfont\MakeUppercase{\thesistitle}\par}
\vspace{1.5cm}
\includegraphics[width=4cm]{logo-isi.png}\par
\vspace{1.5cm}
Oleh:\\
\researchername\\
NIM: \researchernim\par
\vspace{1.5cm}
\mdseries
Skripsi ini diajukan kepada Dewan Penguji \facultyname{} \institutionname{} sebagai salah satu syarat untuk mengakhiri jenjang studi Sarjana S-1 dalam bidang \departmentname\par
\vspace{0.5cm}
\academicyear\par
\end{center}
\vspace*{\fill}
\end{spacing}

% Approval Page
\newpage
\begin{spacing}{1.15}
\begin{center}
    HALAMAN PENGESAHAN\par
    \vspace{0.8cm}
\end{center}
\noindent Skripsi berjudul:\par
\begin{center}
    \textbf{\thesistitle}\par
\end{center}
\noindent diajukan oleh \researchername{}, NIM \researchernim, \programstudy, Jurusan \departmentname, \facultyname{} \institutionname, telah dipertanggungjawabkan di depan Tim Penguji Skripsi pada tanggal \makebox[4cm]{\dotfill} dan dinyatakan telah memenuhi syarat untuk diterima.\par
\vspace{1.0cm}
\begin{center}
    \cityname, \submissiondate\par
    Disetujui,\par
    \vspace{0.4cm}
    
    \begin{minipage}[t]{0.49\textwidth}
        \raggedright
        Penguji I\par
        \vspace{2.0cm}
        \underline{\textbf{\examinerone}}\par
        NIP \examineronenip
    \end{minipage}
    \hfill
    \begin{minipage}[t]{0.49\textwidth}
        \raggedleft
        Penguji II\par
        \vspace{2.0cm}
        \underline{\textbf{\examinertwo}}\par
        NIP \examinertwonip
    \end{minipage}\par
    
    \vspace{1.2cm}
    
    \begin{minipage}[t]{0.49\textwidth}
        \raggedright
        Mengetahui,\par
        Koordinator Program Studi Musik\par
        \vspace{2.0cm}
        \underline{\textbf{\advisoracademic}}\par
        NIP \advisoracademicnip
    \end{minipage}
    \hfill
    \begin{minipage}[t]{0.49\textwidth}
        \raggedleft
        ~\par
        Pembimbing Skripsi\par
        \vspace{2.0cm}
        \underline{\textbf{\advisorthesis}}\par
        NIP \advisorthesisnip
    \end{minipage}
\end{center}
\end{spacing}

% Statement Page
\newpage
\begin{spacing}{1.15}
\begin{center}
    \textbf{HALAMAN PERNYATAAN}
\end{center}
\vspace{0.5cm}

Saya menyatakan bahwa skripsi ini bebas dari unsur plagiasi dan belum pernah diajukan untuk memperoleh derajat akademik di suatu perguruan tinggi. Skripsi ini tidak memuat karya atau pendapat yang pernah ditulis atau diterbitkan oleh orang lain dan/atau oleh saya sendiri sebelumnya, kecuali yang secara tertulis diacu dalam naskah dan tercantum dalam daftar pustaka. Pernyataan ini saya buat dengan sesungguhnya. Saya bersedia menerima sanksi sesuai dengan ketentuan yang berlaku apabila di kemudian hari ditemukan bukti bahwa pernyataan ini tidak benar.\par

\vspace{1.0cm}
\begin{flushright}
    \begin{minipage}{0.45\textwidth}
        \centering
        \cityname, \makebox[3cm]{\dotfill}\par
        Yang menyatakan,\par
        \vspace{2.0cm}
        \researchername\par
        NIM \researchernim
    \end{minipage}
\end{flushright}
\end{spacing}

% Abstract
\newpage
\begin{spacing}{1.15}
\begin{center}
    \textbf{ABSTRAK}
\end{center}

Penelitian ini mengevaluasi akurasi \textit{output} model kecerdasan buatan (AI) generatif musik simbolik terhadap kaidah harmoni fungsional Barat, dengan membandingkan tiga pendekatan arsitektur yang berbeda: \textit{Recurrent} (LSTM, diwakili DeepBach), Konvolusional (CNN, diwakili Coconet), dan \textit{Attention-based} (\textit{Transformer}, diwakili NotaGen). Dalam ranah ilmu komputer dan \textit{Music Information Retrieval} (MIR), kaidah harmoni fungsional Barat diposisikan sebagai \textit{deterministic constraint}, yaitu batasan biner yang dapat diverifikasi secara algoritmik, bukan sebagai standar estetika seni. Posisi ini dipakai untuk menguji sejauh mana model generatif berbasis probabilitas statistik mampu menangkap aturan yang dilarang secara absolut, bukan sekadar aturan yang jarang muncul dalam data latih. Evaluasi dilakukan secara algoritmik menggunakan pustaka music21 dengan tiga kriteria yang diformalkan dari teori harmoni fungsional Gustav Strube (1928): larangan kuint sejajar, larangan oktaf sejajar, dan resolusi nada penuntun. Hasil penelitian diharapkan menunjukkan sejauh mana pergeseran paradigma arsitektur AI, dari \textit{Recurrent} dan Konvolusional generasi 2017 menuju \textit{Attention-based} (\textit{Transformer}) generasi 2025, memengaruhi \textit{constraint adherence} (tingkat kepatuhan \textit{output} model terhadap kaidah harmoni formal) model tersebut.\par

\textbf{Kata Kunci:} Musik Simbolik, \textit{Constraint Adherence}, Harmoni Fungsional, Gustav Strube, LSTM, CNN, \textit{Transformer}
\end{spacing}

% Abstract (English)
\newpage
\begin{spacing}{1.15}
\begin{center}
    \textbf{ABSTRACT}
\end{center}

This study evaluates the constraint adherence of symbolic music generative Artificial Intelligence (AI) models against Western functional harmony rules, comparing three distinct architectural paradigms: Recurrent (LSTM, represented by DeepBach), Convolutional (CNN, represented by Coconet), and Attention-based (\textit{Transformer}, represented by NotaGen). Within computer science and Music Information Retrieval (MIR), Western functional harmony rules are reframed as deterministic constraints, binary in nature and algorithmically verifiable, rather than aesthetic dogma. This framing tests the extent to which probability-based generative models capture rules that are absolutely forbidden, as opposed to rules that are merely rare in training data. Evaluation is performed algorithmically using the music21 library, with three criteria formalized from Gustav Strube's (1928) theory of functional harmony: parallel fifths, parallel octaves, and leading tone resolution. Results are expected to show the extent to which the architectural paradigm shift, from Recurrent and Convolutional models in 2017 to Attention-based (\textit{Transformer}) models in 2025, affects constraint adherence.\par

\textbf{Keywords:} Symbolic Music, Constraint Adherence, Functional Harmony, Gustav Strube, LSTM, CNN, \textit{Transformer}
\end{spacing}
\newpage
